\subsection{Maxwell's equation}
Hydrocarbon materials such that one used in this project can be reasonably assumed to be isotropic and homogeneous. This simplifies the Maxwell's equations in matter because it allows two vectors that are difficult
to compute, namely the polarization vector $\mathbf{P}$ and the magnetization vector $\mathbf{M}$, to be parallel to the electric and magnetic fields. Moreover, $\mu_r \approx 1$ is assumed since hydrocarbon are
non-magnetic materials. This leads to the following form of Maxwell's equations in matter:

\begin{subequations}
    \begin{align}
        \nabla \cdot \mathbf{E} &= \frac{\rho_f}{\epsilon_r \epsilon_0} \label{eqn:gauss} \\
        \nabla \cdot \mathbf{B} &= 0 \label{eqn:gauss_magnetic} \\
        \nabla \times \mathbf{E} &= -\frac{\partial \mathbf{B}}{\partial t} \label{eqn:faraday} \\
        \nabla \times \mathbf{B} &= \mu_0 \mathbf{j}_f + \mu_0 \epsilon_r \epsilon_0 \frac{\partial \mathbf{E}}{\partial t} \label{eqn:ampere}
    \end{align}
    \label{eqn:maxwell}
\end{subequations}

In plasma, the free charge density $\rho_f$ and current density $\mathbf{j}_f$ can be easily computed using electron and ion information:
\begin{subequations}
    \begin{align}
        \rho_f &\equiv e(N_i - N_e) \\
        \mathbf{j}_f &\equiv e(N_i \mathbf{u}_i - N_e \mathbf{u}_e)
    \end{align}
\end{subequations}

% The relative permittivity $\varepsilon_r$ is related to the refractive index $n$ through a simple relation:
% \begin{equation}
%     \varepsilon_r = n^2,
% \end{equation}
% due to the fact that light in matter must travel at the speed $v = c/n$.

Though Equation \ref{eqn:maxwell} explicitly describes four vector equations, it can be easily shown that taking the divergence of Equations \ref{eqn:faraday} and \ref{eqn:ampere} will lead to Equations \ref{eqn:gauss_magnetic}
and \ref{eqn:gauss}, respectively. Therefore, only Equations \ref{eqn:faraday} and \ref{eqn:ampere} are necessary to model the $\mathbf{E}$ and $\mathbf{B}$ fields.

\subsection{Lorentz force}
The Lorentz force acts upon any charge particle has the form
\begin{equation}
    \mathbf{f}_j = \pm e \left(\mathbf{E} + \mathbf{u}_j \times \mathbf{B} \right),
    \label{eqn:lorentz_force}
\end{equation}
here the subscript $j$ can be either $e$ or $i$. Since a force is the rate of change of a momentum, Equation \ref{eqn:lorentz_force} can be expressed as 

\begin{equation}
    m_j \frac{d N_j\mathbf{u}_j}{dt} = \pm eN_j \left(\mathbf{E} + \mathbf{u}_j \times \mathbf{B} \right).
    \label{eqn:lorentz_momentum}
\end{equation}

Since electrons are extremely lightweight, its acceleration due to the Lorentz force easily dominates any diffusive or gravitational effects, so they are safe to neglect. The case for ions is a little bit more nuanced;
however, the Lorentz force is still expected to be the main momentum contributor.

\subsection{Energy transport by electromagnetic waves}
Electromagnetic waves themselves carry energy, and therefore capable of doing work on a charged system. According to Equation \ref{eqn:lorentz_momentum}, the magnetic field does no work on charged particles as its
force component is always perpendicular to the charge velocity. The rate of work done by the electric field is then
\begin{equation}
    \frac{\partial N_j e_j}{\partial t} = m_j \frac{d N_j\mathbf{u}_j}{dt} \cdot \mathbf{u}_j = \pm eN_j \mathbf{u}_j \cdot \mathbf{E},
\end{equation}
and should be equal to the rate of energy gained by the particles.